% !TeX program = xelatex
% UTF-8. Standalone XeLaTeX/XePersian document.
% Define \ElectrolyteDisordersInNotebook before \input in a XePersian parent.
% Source: Google Drive folder 16ldoYts9x3By1XTZGRvt89dDz9PhHG2p.
% Original photographs are preserved unchanged, in ascending filename order.
% pic01.jpg = 20261016_205058238.jpg | 1GfnhQheNzsTqMVuuoykFyCtAZtGt_ETd
% pic02.jpg = 20261016_205104925.jpg | 12Bv3TOEgkWCqs-9V-HZeSCE2X-A375Jd
% pic03.jpg = 20261016_205109133.jpg | 1vo5mpX8XhO3rpFaAGdaVbs_ucCz9i6Wl
% pic04.jpg = 20261016_205116221.jpg | 1gE9mRGB1w_rXAijgY_IV5ZoUzBSrGETX
% pic05.jpg = 20261016_205122733.jpg | 1gB0YABDtRnmb2sonsNjGeXpVh5CIhGUI
% pic06.jpg = 20261016_205129947.jpg | 1_1V6np8sGuunKPiLklHIUL208a0T_d4s
% pic07.jpg = 20261016_205135102.jpg | 1fH0eHVyYAjvGLXVhvbOa6GHEZ3hmJAvd
% pic08.jpg = 20261016_205140876.jpg | 1vqA2TMlPQ2XImPBXNtG4zxXdSgK3M6lN
% pic09.jpg = 20261016_205145809.jpg | 1008WQ1MfbIvUafOTAvr0tlWQxlhetoRQ
% pic10.jpg = 20261016_205151134.jpg | 1RjyLBbt25wcL4Z4wDU9aAzAK7Wv8h4H3
% pic11.jpg = 20261016_205156510.jpg | 1iInTcWAv73ChwBgr3sCkoFL-JL3ynXg0
% pic12.jpg = 20261016_205200776.jpg | 1ue0FKeyfmICNRN2D4tEZqpKBWni_NiyJ
\ifdefined\ElectrolyteDisordersInNotebook
  \let\EDfinish\relax
\else
  \def\EDfinish{\end{document}}
  \documentclass[12pt,a4paper]{article}
  \usepackage[margin=22mm]{geometry}
  \usepackage{amsmath,amssymb}
  \usepackage{graphicx}
  \usepackage{enumitem}
  \usepackage{xepersian}
  \IfFontExistsTF{Amiri}{\settextfont{Amiri}}{\settextfont{DejaVu Sans}}
  \setlatintextfont{DejaVu Serif}
  \setlength{\parindent}{0pt}
  \setlength{\parskip}{5pt}
  \setlength{\emergencystretch}{3em}
  \begin{document}
\fi
\providecommand{\EDm}[1]{\lr{\(#1\)}}
\providecommand{\EDunclear}[1]{\textbf{[خوانش نامطمئن: #1]}}
\providecommand{\EDpage}[1]{\clearpage\section*{تصویر #1 ـ صفحهٔ دست‌نویس #1}}
% Printed stationery on every sheet: Subject: / Year: / Month: / Date:
% Printed line numbers on every sheet: 1, 2, 3, 4, 5, 6, 7, 8, 9.
% Empty stationery fields are recorded below rather than filled in.
% The transcription preserves source statements, even when medically unusual.
% Arrows and braces are expressed as causal sequences, headings and lists.
\begin{center}
{\Large\bfseries آب و الکترولیت کودکان}
\end{center}
\section*{تصویر ۱ ـ صفحهٔ دست‌نویس ۱}
\lr{Subject:} آب و الکترولیت کودکان.\par
\lr{Year: \quad Month: \quad Date:} خالی.\par
در تصاویر ۲ تا ۱۲، فیلد \lr{Subject:} نیز خالی است. شماره‌های چاپی سطرهای هر تصویر: \lr{1, 2, 3, 4, 5, 6, 7, 8, 9}.

\subsection*{\lr{Total Body Water (TBW)}}
\begin{itemize}
\item در نوزاد ترم ۷۵\% است؛ در نوزاد پره‌ترم بیشتر است.
\item در یک‌سالگی: ۶۰\% می‌رسد؛ تا بلوغ همین می‌ماند.
\item پس از بلوغ: در \lr{M} همان ۶۰\% و در \lr{F}، ۵۰\% می‌شود.
\item در یک‌سالگی نسبت \EDm{\dfrac{ICF}{ECF}} مشابه بزرگسالان می‌شود.
\item \lr{ECF} تقریباً ۲۰--۲۵\% از \lr{BW} است.
\item \lr{ICF} تقریباً ۳۰--۴۰\% از \lr{BW} است.
\item \lr{ECF}: پلاسما معادل ۵\% \lr{BW}؛ مایع بینابینی معادل ۱۵\% \lr{BW}.
\item حجم خون با فرض \EDm{HCT=40\%}، حدود ۸\% \lr{BW} است.
\end{itemize}
\subsection*{اسمولالیته}
در حالت عادی اسمولالیتهٔ \lr{ICF} و \lr{ECF} برابر در نظر گرفته می‌شود.
\begin{latin}
Plasma Osmolality: N/L = 285--295 mOsm/kg.
\end{latin}
در آزمایشگاه با استفاده از روش فریزینگ پوینت دپرشن اندازه‌گیری می‌شود.
\begin{latin}
Estimation:
\[
\mathrm{Osmolality}=2\times[\mathrm{Na}]+\frac{[\mathrm{Glu}]}{18}+\frac{[\mathrm{BUN}]}{2.8}
\]
\(\mathrm{mg/dL}\longrightarrow\mathrm{mmol/L}\)
\end{latin}
برای در نظر گرفتن آنیون‌ها، سدیم را در ۲ ضرب می‌کنیم (کلر و بی‌کربنات).

\lr{BUN} بدون محدودیت در عرض غشا حرکت می‌کند و به‌جز در \lr{DES}، از نظر اسمولالیته اهمیت بالینی ندارد.

اتانول نیز بدون محدودیت از غشا عبور می‌کند و یک اسمول غیرمؤثر است.
\begin{latin}
\[
\mathrm{Effective\ Osmolality}=2\times[\mathrm{Na}]+\frac{[\mathrm{Glu}]}{18}
\]
\end{latin}

\EDpage{۲}
\subsection*{هایپرگلاسمی و سدیم تصحیح‌شده}
در هایپرگلاسمی اسمولالیتهٔ مؤثر \lr{ECF} بالا می‌رود؛ آب از \lr{ICF} وارد \lr{ECF} می‌شود؛ یون \lr{Na} در \lr{ECF} رقیق می‌شود؛ علی‌رغم بالا رفتن اسمولالیته، هایپوناترمی رخ می‌دهد.
\begin{latin}
\[
[\mathrm{Na}]_{\mathrm{Corrected}}=[\mathrm{Na}]_{\mathrm{Measured}}+
1.6\times\frac{[\mathrm{Glu}]-100\,\mathrm{mg/dL}}{100}
\]
\end{latin}
سدیم اصلاح‌شده معیار معتبرتری برای ارزیابی نسبت \lr{Na} کل بدن به \lr{TBW} است.

در حالت عادی اختلاف اسمولالیتهٔ محاسبه‌شده و اندازه‌گیری‌شده کمتر از \lr{10 mOsm/kg} است.

\lr{Osmolal Gap}: به معنای وجود اسمول‌های اندازه‌گیری‌نشده در خون است؛ اتانول، متانول، اتیلن گلیکول، سایر قندها، سوربیتول، مانیتول.

\subsection*{سودوهایپوناترمی}
\begin{itemize}
\item در افرادی که لیپید یا پروتئین بالا دارند، آب در سرم کم است، یا پروتئین جایگزین شده است.
\item برخی دستگاه‌ها غلظت \lr{Na} در کل سرم را اندازه می‌گیرند، نه در آب سرم؛ نشان می‌دهد ولی چیزی که مهم است غلظت \lr{Na} در آب سرم است!
\item علی‌رغم \lr{Na} کم، تست آزمایشگاهی نشان می‌دهد اسمولالیتهٔ اندازه‌گیری‌شده \lr{N/L} است.
\item در دستگاه‌های \lr{ABG} این مشکل رخ نمی‌دهد.
\end{itemize}
کاربرد محاسبهٔ اسمولالیته: تشخیص وجود اسمول‌های اندازه‌گیری‌نشده؛ تأیید حقیقی بودن هایپوناترمی.

شایع‌ترین علت افزایش اسمولالیته در کودکان، دهیدریشن است.

\EDpage{۳}
\subsection*{روش‌های تنظیم اسمولالیته و حجم در کودکان}
\begin{itemize}
\item حفظ اسمولالیتهٔ نرمال وابسته به کنترل بالانس مایع (آب) است.
\item کنترل وضعیت حجمی وابسته به کنترل وضعیت سدیم است.
\end{itemize}
حفظ حجم داخل عروقی برای بدن مهم‌تر از حفظ اسمولالیته است.

\subsection*{تنظیم اسمولالیته}
\begin{itemize}
\item تشنگی و دفع مایع از کلیه، مکانیسم‌های اصلی هستند.
\item اسمورسپتورها در هیپوتالاموس.
\item اسمولالیتهٔ پلاسما: درک می‌کنند.
\item \lr{Effective Osmolality}: ترشح \lr{ADH} از نورون‌های هستهٔ سوپرااپتیک و پاراونتریکولار در هیپوتالاموس.
\item \lr{ADH} به گیرندهٔ \EDm{V_2} در \lr{CD} نفرون؛ کانال‌های \lr{AQP2} وارد غشای سلول‌ها می‌شوند؛ آب بازجذب.
\item احساس تشنگی: در اثر افزایش اسمولالیتهٔ پلاسما یا \lr{Volume Depletion} می‌تواند ایجاد شود.
\end{itemize}
در \lr{SIADH} علی‌رغم اسمولالیتهٔ پایین پلاسما، \lr{ADH} همچنان تولید می‌شود.

افت \lr{GFR}: توانایی کلیه در بازجذب آب را کم می‌کند چون آب کمتری به دیستال \lr{CD} می‌رسد؛ ولی برای اختلال در دفع آب زیادی باید نارسایی کلیه خیلی شدید باشد.

حداقل اسمولالیتهٔ ادرار: \lr{30--50 mOsm/kg}؛ حداکثر میزانی که کلیه می‌تواند آب دفع کند؛ از این بیشتر منجر به مسمومیت با آب می‌شود.

یک شرایط خاص که در آن \lr{ADH} به جای پاسخ به اسمولالیتهٔ مؤثر، به حجم آب پاسخ می‌دهد: یعنی در واقع \lr{Effective Osm} پایین است ولی ما \lr{Volume Depletion} داریم؛ در پاسخ \lr{ADH} ترشح می‌شود. این یک پاسخ حیاتی و مناسب است؛ در نوزادان و شیرخواران مهم است.

\EDpage{۴}
حداکثر اسمولالیتهٔ ادرار: \lr{1200 mOsm/kg} است.

حداقل حجم ادرار، پس از یک \lr{AKI} یا پس از رفع یک انسداد، افزایش می‌یابد تا \lr{BUN} اضافه را دفع کند.

یک نوزاد توانایی کمتری در تغلیظ ادرار دارد؛ خصوصاً نوزاد پره‌مچور؛ استعداد بالا برای ابتلا به دهیدریشن هایپرناترمیک.

افرادی که حجم مایع بسیار زیاد مصرف می‌کنند: مانند پلی‌دیپسی سایکوژنیک.
\begin{itemize}
\item اسمولالیتهٔ مدولای کلیه به طور طبیعی بالاست و باعث می‌شود کلیه بتواند ادرار را تغلیظ کند؛ در این افراد رقیق می‌شود.
\item وقتی به طور ناگهانی مصرف مایعات کم می‌شود، توانایی کلیه برای تغلیظ ادرار کاهش می‌یابد؛ این مشکل معمولاً پس از چند روز رفع می‌شود.
\item این مشکل در شروع درمان \lr{DI} سنترال نیز رخ می‌دهد؛ زمانی که درمان با دسموپرسین را شروع می‌کنیم کلیه نمی‌تواند ادرار را غلیظ کند.
\end{itemize}
\subsection*{تنظیم حجم}
\begin{itemize}
\item وجود میزان کافی سدیم در بدن برای حفظ حجم داخل عروقی ضروری است.
\item کلر هم به اندازهٔ سدیم اهمیت دارد ولی برای ساده‌سازی ما فقط به سدیم اشاره می‌کنیم.
\item برای مثال: آلکالوز متابولیک \lr{Volume depletion} که علتش کلر است.
\item تنظیم غلظت سدیم بر عهدهٔ کلیه است. ما اغلب کنترلی بر مصرف و ورود سدیم به بدن نداریم؛ غیر \lr{Salt Craving} که در کودکان با \lr{Loss} مزمن کلیوی نمک رخ می‌دهد.
\item کلیه از طریق تغییر در مقداری از \lr{Na} که بازجذب می‌شود، غلظت آن را تنظیم می‌کند.
\item در حالت عادی کمتر از ۱\% سدیم فیلترشده از بدن دفع می‌شود.
\item در فرد سالم: ورودی سدیم با دفع ادراری + دفع خارج کلیوی سدیم برابر است.
\end{itemize}

\EDpage{۵}
در مواقع نیاز دفع کلیوی سدیم می‌تواند عملاً به صفر برسد و یا برعکس، شدیداً بالا برود.

مهم‌ترین عامل تعیین‌کنندهٔ میزان دفع سدیم، وضعیت حجمی کودک است.
\begin{itemize}
\item \lr{EIV}: حجم مؤثر داخل عروقی.
\item در نارسایی قلب با اینکه آب بدن زیاد است ولی \lr{EIV} پایین است.
\item به همین دلیل در \lr{HF}، \lr{Na} احتباس پیدا می‌کند.
\end{itemize}
سیستم \lr{RAAS} یک تنظیم‌کنندهٔ مهم دفع سدیم است.
\begin{itemize}
\item در پاسخ به کاهش \lr{EIV}، رنین از جاگستاگلومرولار ترشح می‌شود.
\item تحریک اصلی: کاهش پرفیوژن کلیه در شریان \lr{Afferent} است.
\item کاهش دلیوری سدیم به دیستال نفرون و آگونیست‌های \EDm{\beta_1} که در پاسخ به کاهش \lr{EIV}، رنین را بالا می‌برند.
\item در نهایت آنژیوتانسین \lr{II} با اثر بر \lr{PCT} کلیه: افزایش بازجذب \lr{Na}.
\item تحریک آدرنال: افزایش ترشح آلدوسترون.
\item آلدوسترون با اثر بر \lr{late DCT} و \lr{CD}: افزایش بازجذب \lr{Na}؛ افزایش ترشح \lr{K}؛ افزایش دفع ادراری.
\item آنژیوتانسین \lr{II} همچنین یک وازوکانستریکتور عروقی نیز می‌باشد.
\end{itemize}
اگر حجم داخل عروقی مؤثر (\lr{EIV}) بالا برود: تحریک آزادسازی \lr{ANP} از دهلیزها (آتریال ناتریورتیک پپتاید).
\begin{itemize}
\item \lr{ANP}: افزایش \lr{GFR}.
\item مهار بازجذب \lr{Na} در بخش مدولاری \lr{CD}.
\end{itemize}
زمانی که دریافت \lr{Na} از دفع آن بیشتر شود، \lr{Volume Overload} رخ می‌دهد؛ کوشینگ، \lr{AKI}، نوزادان پره‌مچور.

\EDpage{۶}
در برخی بیماری‌ها تنظیم صحیح \lr{Na} مختل می‌شود.
\begin{itemize}
\item بیمار با آلدوسترون بالا: هایپرآلدوسترونیسم اولیه؛ استنوز شریان رنال؛ به علت رنین بالا.
\item گلومرولونفریت حاد حتی بدون افت شدید قابل توجه \lr{GFR}، به دلیل اختلال مکانیسم‌های داخل کلیوی تنظیم ترشح \lr{Na}: احتباس سدیم و \lr{Overload} مایع رخ می‌دهد.
\end{itemize}
احتباس \lr{Na} توسط کلیه به دلیل \lr{Volume depletion} یک پاسخ درست است ولی منجر به افزایش شدید \lr{Na} توتال بدن می‌شود.
\begin{itemize}
\item مثال: نارسایی قلب، نارسایی کلیه، سندرم نفروتیک، سایر علل هایپوآلبومینمی.
\item در نتیجه: ادم، افزایش فشار خون رخ می‌دهد.
\end{itemize}
زمانی که دفع سدیم بدن بیشتر از دریافت سدیم باشد، \lr{Volume Depletion} رخ می‌دهد.
\begin{itemize}
\item شایع‌ترین مثال گاستروانتریت است.
\item سوءتغذیه؛ در بیماران \lr{CF} به علت دفع زیاد سدیم در عرق.
\end{itemize}
دریافت ناکافی \lr{Na} ناشایع است مگر در \lr{Neglect}، \lr{Famine}، مایع‌درمانی اشتباه.

\lr{Urinary Na wasting} در برخی بیماری‌های کلیه رخ می‌دهد:
\begin{itemize}
\item رنال دیسپلازی.
\item اختلالات توبولار: سندرم بارتر.
\item نوزادان خصوصاً پره‌مچور، مقداری این مشکل را دارند.
\item علل ایاتروژنیک: درمان با دیورتیک.
\item کمبود آلدوسترون: هایپرپلازی \EDunclear{همراه با تولد} آدرنال (کمبود ۲۱ ـ هیدروکسیلاز).
\end{itemize}

\EDpage{۷}
اختلالات ایزولهٔ تنظیم آب در بدن نیز می‌توانند وضعیت حجم و \lr{Na} را تحت تأثیر قرار دهند.
\begin{itemize}
\item در \lr{Isolated water loss} مانند \lr{DI}، حجم \lr{ICF} بیشتر تحت تأثیر قرار می‌گیرد.
\item به همین علت دهیدریشن هایپرناترمیک اثر کمتری بر حجم پلاسما در مقایسه با سایر علل دهیدریشن دارد؛ چون اکثر مایع مربوط به \lr{ICF} است.
\end{itemize}
از دست رفتن حجم اگر زیاد باشد، حتی اگر سدیم توتال نرمال باشد، منجر به افزایش بازجذب سدیم در کلیه می‌شود.

در مسمومیت حاد با آب، \lr{SIADH}، \lr{TBW} بالا می‌رود ولی اکثر آن در \lr{ICS} است. اگرچه که در حدی حجم داخل عروقی نیز بیشتر است که منجر به افزایش ترشح \lr{Na} در کلیه می‌شود با اینکه سدیم توتال بدن نرمال است.

کودکان با مسمومیت حاد با آب و یا \lr{SIADH}: سدیم ادرار بالا علی‌رغم هایپوناترمی.

\subsection*{سدیم}
\begin{itemize}
\item ۳\% از سدیم در داخل سلول است. حدود ۴۰\% آن در استخوان است.
\item بقیه در فضای داخل عروقی و بین‌سلولی.
\item سدیم داخل سلول در حد \lr{10 mEq/L} توسط پمپ \lr{Na-K-ATPase} حفظ می‌شود.
\item سدیم شیر مادر \lr{7 mEq/L}، سدیم فرمولا \lr{7--13 mEq/L} به ازای فرمولای \lr{20 Cal/oz} است.
\item سدیم در تمام طول \lr{GI} جذب می‌شود. مینرالوکورتیکوئیدها جذب سدیم از \lr{GI} را بالا می‌برند، ولی اثر بالینی قابل توجهی ندارد.
\item وجود گلوکز جذب سدیم را افزایش می‌دهد (وجود کوترنسپورتر)؛ به این علت در \lr{ORS} سدیم و گلوکز وجود دارد.
\end{itemize}

\EDpage{۸}
\subsection*{دفع سدیم و تنظیم غلظت آن}
\begin{itemize}
\item دفع سدیم: پوست، مدفوع، ادرار.
\item دفع سدیم در مدفوع بی‌اهمیت است مگر آنکه اسهال وجود داشته باشد.
\item سدیم عرق \lr{5--40 mEq/L} است. در کودک \lr{CF} بالاتر.
\item در کودک با آلدوسترون دفیشنسی و سودوهایپرآلدوسترونیسم بالاتر است.
\end{itemize}
غلظت سدیم وابسته به بالانس سدیم نیست بلکه وابسته به بالانس آب است.

افزایش سدیم؛ افزایش اسمولالیتهٔ پلاسما؛ افزایش تشنگی و افزایش ترشح \lr{ADH}؛ افزایش حجم آب و در نهایت کاهش غلظت سدیم.

بالانس آب به طور معمول توسط اسمولالیتهٔ پلاسما کنترل می‌شود ولی کاهش حجم آب نیز تشنگی را تحریک می‌کند و ترشح \lr{ADH} را بالا می‌برد.

در کلیه، ترشح سدیم توسط اسمولالیتهٔ پلاسما کنترل نمی‌شود بلکه \lr{EIV} ترشح یا احتباس سدیم را کنترل می‌کند؛ توسط محور \lr{RAAS} و مکانیسم‌های داخل نفرون.

\subsection*{هایپوناترمی در کودکان}
\begin{itemize}
\item در کودکان بستری بسیار شایع است.
\item در سه حالت \lr{low}، نرمال، \lr{high} از \lr{TBW} ممکن است رخ دهد.
\item و در سه حالت \lr{low}، \lr{N/L}، \lr{high} در \lr{Total Na} ممکن است رخ دهد.
\end{itemize}
\subsection*{اتیولوژی و پاتوفیزیولوژی}
\subsubsection*{سودوهایپوناترمی}
وقتی سطح پروتئین‌ها یا چربی‌ها \lr{Very high} باشد: \lr{MM}، تزریق \lr{IVIG}، هایپرلیپیدمی (\EDm{\uparrow TG}، \EDm{\uparrow Chol}).

بهتر است در این شرایط تست با روش زیر انجام شود:
\begin{latin}
Direct ion-selective electrode
\end{latin}

\EDpage{۹}
افتراق سودو از ترو: در هایپوناترمی واقعی، اسمولالیتهٔ اندازه‌گیری‌شده پایین است.

هایپراسمولالیته (مثلاً در هایپرگلاسمی) منجر به هایپوناترمی می‌شود چون آب از \lr{ICF} به \lr{ECF} می‌آید.
این بیماران علائم هایپوناترمی ندارند چون علائم عصبی هایپوناترمی به علت هایپواسمولالیته رخ می‌دهد که در این بیماران وجود ندارد.

مانیتول و سوکروز که در فرآوردهٔ \lr{IVIG} هستند می‌توانند با مکانیسم هایپراسمولالیته، هایپوناترمی ایجاد کنند.

تقسیم‌بندی هایپوناترمی بر اساس وضعیت حجمی بیمار است.
\subsection*{۱) هایپوناترمی هایپوولمیک}
کودک \lr{Na} از دست داده است. بالانس آب می‌تواند مثبت یا منفی باشد ولی دفع \lr{Na} بیشتر از از دست دادن آب بوده است.

پاتوژنز هایپوناترمی اغلب: ترکیبی از \lr{Na loss} و \lr{water retention} برای جبران \lr{Volume depletion}.

بیمار اغلب به طور پاتولوژیک دفع افزایش‌یافتهٔ \lr{fluid} دارد و این مایع حاوی \lr{Na} است.
\begin{itemize}
\item اغلب مایعات، غلظت سدیم کمتر از پلاسما دارند اما شدت آن است که مایعی که جایگزین می‌کنیم ممکن است حتی سدیم کمتری داشته باشد.
\item غلظت سدیم در اسهال ویروسی \lr{50 mEq/L} است ولی در شیرخوار \lr{7--10 mEq/L}.
\end{itemize}
مکانیسم معمول بدن برای مقابله با هایپوناترمی، دفع آب است. این مکانیسم در بیماری که اسهال و دچار \lr{Volume Depletion} شده است، به مشکل می‌خورد.

اسهال شایع‌ترین علت هایپوناترمی هایپوولمیک است در کودکان.

استفراغ اگر با مایع هایپوتونیک وریدی یا خوراکی جایگزین شود: هایپوناترمی.

\EDpage{۱۰}
اکثر بیماران با استفراغ یا \lr{Na} نرمال دارند یا هایپرناترمی هستند.

دریافت سرم هایپوتونیک: هایپوناترمی.

از دست دادن مایع در فضای سوم: مایع تقریباً سرم ایزوتونیک است و منجر به کاهش \lr{EIV} می‌شود؛ \EDm{\uparrow ADH}؛ \lr{Water retention}؛ هایپوناترمی.

در بیماری‌هایی که منجر به از دست رفتن غیرکلیوی \lr{Na} می‌شوند، سدیم ادرار باید کمتر از \lr{10 mEq/L} باشد (به علت بازجذب کلیه). استثنا بیماری‌هایی هستند که منجر به از دست رفتن کلیوی و غیرکلیوی همزمان \lr{Na} می‌شوند: نارسایی آدرنال و سودوهایپوآلدوسترونیسم.

\subsection*{هایپوناترمی هایپوولمیک}
\begin{itemize}
\item \lr{Extrarenal loss}: اسهال، استفراغ، سوختگی، تعریق، فضای سوم.
\item \lr{Renal loss}: دیورتیک لوپ و تیازید، دیورز اسموتیک، دیورز پست‌ابستراکتیو، فاز پلی‌اوریک \lr{ATN}، سودوهایپوآلدوسترونیسم، آدیسون، فرم \lr{AR} کلیهٔ پلی‌کیستیک (\lr{PCKD})، \lr{RTA} پروگزیمال (تیپ \lr{II})*.
\end{itemize}
\subsection*{هایپوناترمی یوولمیک}
\lr{SIADH}، کمبود گلوکوکورتیکوئید، هایپوتیروئیدی، داروهای ضدافسردگی، مسمومیت آب.

\subsection*{هایپوناترمی هایپرولمیک}
\lr{HF}، سیروز، سندرم نفروتیک، \lr{CKD}، هایپوآلبومینمی در زمینهٔ \lr{Loss} گوارشی و \EDunclear{ادامهٔ عبارت در حاشیه کامل نیست}.

* \lr{RTA} پروگزیمال در اکثر موارد علت ژنتیکی ندارد بلکه جزئی از سندرم فانکونی است که می‌تواند اتیولوژی‌های مختلفی داشته باشد.

در برخی موارد ممکن است سدیم ادرار بالاتر از حد مورد انتظار باشد (\EDm{>20}) که در این صورت حتی بدون تغییر در دریافت مایع نیز هایپوناترمی می‌شود.

\EDpage{۱۱}
زمانی می‌گوییم دفع کلیوی در هایپوناترمی نقش داشته که سدیم ادرار \EDm{>20} باشد.

\subsection*{دیورتیک‌ها}
\begin{itemize}
\item به علت دفع کلیوی ممکن است باعث هایپوناترمی بشوند ولی اگر در زمان آزمایش ادرار قطع شده باشند، سطح سدیم ادرار \EDm{<20} خواهد بود.
\item دیورتیک لوپ اجازه نمی‌دهد مدولای کلیه به حداکثر هایپرتونیسیتهٔ خود برسد؛ به همین دلیل کلیه نخواهد توانست ادرار را در حداکثر شدت خود غلیظ یا رقیق کند؛ بیمار دچار هایپوناترمی شدید نمی‌شود.
\item دیورتیک تیازیدی این اثر را ندارد؛ خطر هایپوناترمی شدید وجود دارد.
\end{itemize}
در مرحلهٔ دیورز اسموتیک و پست‌ابستراکتیو، اگر جایگزینی مایع با سدیم هایپوتون باشد بیمار هایپوناترمیک خواهد شد.

\lr{Renal Salt wasting}: در بیماری‌های ارثی کلیه رخ می‌دهد؛ \lr{PCKD}؛ \lr{Cerebral salt wasting} هم به علت پپتیدهای ناتریورتیک.

در سندرم فانکونی: افزایش ترشح \lr{Na} و بی‌کربنات در ادرار؛ گلوکزاوری، آمینواسیداوری، هایپرفسفاتوری.

آلدوسترون برای دفع \lr{K} و اسید و بازجذب سدیم در کلیه مهم است.

\subsection*{هایپوناترمی هایپرولمیک}
اضافه‌مایع و سدیم وجود دارد ولی اضافه‌مایع بیشتر است.
\begin{itemize}
\item سدیم ادرار زیر ۱۰ است.
\item به‌جز نارسایی کلیه که \EDm{U_{Na}} ممکن است بالا باشد.
\end{itemize}

\EDpage{۱۲}
در هایپوناترمی به علت ورود مایع از \lr{ECS} به \lr{ICS}، سلول متورم می‌شود که در اکثر بافت‌ها قابل تحمل است به‌جز مغز چون توسط سطوح جمجمه احاطه شده است.

در سندرم هایپوناترمی:
\begin{itemize}
\item \EDm{\uparrow ICP}؛ اختلال جریان خون مغز؛ ایسکمی؛ هرنی \lr{Brain-stem}؛ آپنه.
\item آنورکسی، تهوع و استفراغ، تشنج، لتارژی، گیجی، آژیتیشن.
\item سردرد، تشنج، کما، کاهش رفلکس‌ها.
\item هایپوترمی، تنفس شین ـ استوک.
\item کرامپ و ضعف عضلات.
\item رابدومیولیز در مسمومیت با آب.
\end{itemize}
نورون‌ها و سلول‌های مغز در موارد مزمن با خارج کردن یون‌ها و برخی مولکول‌ها، با کاهش اسمولالیتهٔ داخل سلول باعث کاهش تورم سلول می‌شوند.

% Facsimiles retain all marks, source spelling and any uncertain handwriting.
\clearpage
\section*{تصاویر اصلی}
\providecommand{\EDoriginal}[2]{%
  \clearpage\subsection*{تصویر #1 ـ اصل دست‌نویس}
  \IfFileExists{images/#2}{%
    \includegraphics[width=\linewidth,height=.86\textheight,keepaspectratio]{images/#2}%
  }{%
    \includegraphics[width=\linewidth,height=.86\textheight,keepaspectratio]{pediatrics/ped-nephrology/electrolyte disorders/images/#2}%
  }\par}
\EDoriginal{۱}{pic01.jpg}
\EDoriginal{۲}{pic02.jpg}
\EDoriginal{۳}{pic03.jpg}
\EDoriginal{۴}{pic04.jpg}
\EDoriginal{۵}{pic05.jpg}
\EDoriginal{۶}{pic06.jpg}
\EDoriginal{۷}{pic07.jpg}
\EDoriginal{۸}{pic08.jpg}
\EDoriginal{۹}{pic09.jpg}
\EDoriginal{۱۰}{pic10.jpg}
\EDoriginal{۱۱}{pic11.jpg}
\EDoriginal{۱۲}{pic12.jpg}
\EDfinish
